% =====================================================
% Chapter 12: Conclusion
% Conclusion framework only. No specific outcome, achieved
% value, or success/failure claim is asserted, since no
% simulation or measurement results are available at the time
% of writing this report.
% =====================================================
\chapter{Conclusion}
\label{ch:conclusion}

This project set out to design, simulate, and build an \projecttitle\ that automatically adjusts its gain to hold the output signal level within a target range across a specified range of microphone input levels, as motivated in Chapter~\ref{ch:introduction} and specified in Chapter~\ref{ch:problem}. The theoretical background, system architecture, and circuit design approach for achieving this were presented in Chapters~\ref{ch:theory}--\ref{ch:circuitdesign}.

\section{Summary of Outcome}

\placeholder{Once simulation (Chapter~\ref{ch:simulation}) and hardware measurement (Chapter~\ref{ch:results}) are complete, summarise here, for each target specification listed in Table~\ref{tab:target-specs}, whether it was met, partially met, or not met, quoting the key measured value against its target. State explicitly whether the assigned minimum success criterion (Chapter~\ref{ch:problem}) was achieved.}

\section{Key Limitation}

\placeholder{State the single most significant limitation actually observed in the completed prototype (see Chapter~\ref{ch:safety} for general candidate limitations to consider, and Chapter~\ref{ch:results} for the discussion of measured discrepancies)}

\section{Closing Remarks}

The methodology presented in this report -- combining a low-noise preamplifier stage with a feedback-based automatic gain control loop -- follows established analog design practice for handling wide-dynamic-range audio sources \cite{ref1,ref2,ref5,ref6}. Completion of the outstanding simulation and hardware verification steps, and population of the placeholders throughout this report with the team's actual data, will allow this conclusion to be finalised with concrete, measured evidence rather than general statements.

%% TODO: Do not add any new quantitative claims to this chapter
%% that are not already supported by results presented in
%% Chapters 6 and 8. The conclusion should summarise, not
%% introduce, results.
